% Part: first-order-logic
% Chapter: models-theories
% Section: theories

\documentclass[../../../include/open-logic-section]{subfiles}

\begin{document}

\olfileid{fol}{mat}{the}

\olsection{一阶理论的例子}

\begin{ex}
语言~$\Lang L_<$ 中的严格线性序理论由以下集合公理化：
\begin{align*}
\{\quad & \lforall[x][\lnot x < x], \\
& \lforall[x][\lforall[y][((x < y \lor y <
    x) \lor x = y)]], \\
& \lforall[x][\lforall[y][\lforall[z][((x < y
      \land y < z) \lif x < z)]]] \quad \}
\end{align*}
它完全刻画了预期的结构：每个严格线性序都是该公理系统的模型，反之，若 $R$ 是集合 $X$ 上的线性序，则满足 $\Domain M = X$ 且 $\Assign{<}{M} = R$ 的结构 $\Struct M$ 就是该理论的模型。
\end{ex}

\begin{ex}
在包含常项符号 $\Obj 1$ 与二元函数符号 $\cdot$ 的语言中，群论由以下公理公理化：
\begin{align*}
& \lforall[x][\eq[(x \cdot \Obj 1)][x]]\\
& \lforall[x][\lforall[y][\lforall[z][\eq[(x \cdot (y \cdot z))][((x
          \cdot y) \cdot z)]]]]\\
& \lforall[x][\lexists[y][\eq[(x \cdot y)][\Obj 1]]]
\end{align*}
\end{ex}

\begin{ex}
皮亚诺算术理论由算术语言~$\Lang L_A$ 中的以下语句公理化：
\begin{align*}
& \lforall[x][\lforall[y][(\eq[x'][y'] \lif \eq[x][y])]]\\
& \lforall[x][\eq/[\Obj 0][x']]\\
& \lforall[x][\eq[(x + \Obj 0)][x]]\\
& \lforall[x][\lforall[y][\eq[(x + y')][(x + y)']]]\\
& \lforall[x][\eq[(x \times \Obj 0)][\Obj 0]]\\
& \lforall[x][\lforall[y][\eq[(x \times y')][((x \times y) + x)]]]\\
& \lforall[x][\lforall[y][(x < y \liff \lexists[z][\eq[(z' + x)][y])]]]\\
\intertext{再加上所有形如}
& (!A(\Obj 0) \land \lforall[x][(!A(x) \lif !A(x'))]) \lif \lforall[x][!A(x)]
\end{align*}
由于后一种形式的语句有无穷多条，该公理系统是无穷的。
后一种形式被称为\emph{归纳模式}。（实际上，归纳模式比我们这里呈现的要稍微复杂一些。）

最后一个公理是~$<$ 的\emph{显式定义}。
\end{ex}

\begin{ex}
纯集合理论在数学基础（以及数学哲学）中扮演着重要角色。
如果一个集合的所有元素也是纯集合，则该集合是\emph{纯的}。
因此空集算作纯集合，但如果一个集合包含不是集合的元素，它就不是纯的。
所以，纯集合就是仅由空集生成、不含任何“本元”（即本身不是集合的对象）的集合。

以下可以被视为纯集合理论的一个公理系统：
\begin{align*}
& \lexists[x][\lnot \lexists[y][y \in x]]\\
& \lforall[x][\lforall[y][(\lforall[z](z \in x \liff z \in y) \lif 
\eq[x][y])]]\\
& \lforall[x][\lforall[y][\lexists[z][\lforall[u][(u \in z \liff
(\eq[u][x] \lor \eq[u][y]))]]]]\\
& \lforall[x][\lexists[y][\lforall[z][(z \in y \liff \lexists[u][(z \in
        u \land u \in x)])]]]\\
\intertext{再加上所有形如} &
\lexists[x][\lforall[y][(y \in x \liff !A(y))]]
\end{align*}
第一条公理说存在一个没有元素的集合（即 $\emptyset$ 存在）；第二条说集合具有外延性；第三条说对任意集合 $X$ 和 $Y$，集合 $\{X, Y\}$ 存在；第四条说对任意集合 $X$，集合 $\cup X$ 存在，其中 $\cup X$ 是 $X$ 的所有元素的并集。

最后提到的那些语句统称为\emph{朴素概括模式}。
它本质上断言，对每个 $!A(x)$，集合 $\Setabs{x}{!A(x)}$ 都存在——所以乍一看，这是一个为真的、有用的、甚至可能是必要的公理。
之所以被称为“朴素”，是因为事实证明它使得该理论不可满足：若取 $!A(y)$ 为 $\lnot y \in y$，则会得到语句
\[
\lexists[x][\lforall[y][(y \in x \liff \lnot y \in y)]]
\]
而该语句在任何结构中都不被满足。
\end{ex}

\begin{ex}
在\emph{部分学}领域，\emph{部分关系}是一种基本关系。就像集合论一样，也存在对部分关系的各种（有时是相互冲突的）概念进行公理化的理论。

部分学的语言包含一个单一的二元谓词符号~$\Obj P$，而 $\Atom{\Obj P}{x, y}$ “意指” $x$ 是 $y$ 的一部分。
在此解释下，这种语言的一个结构被称为\emph{部分关系结构}。
当然，并非每个具有单一二元谓词的结构都真正适用这一名称。
为了有可能刻画“部分关系”，$\Assign{\Obj P}{M}$ 必须满足某些条件，我们可以把这些条件作为部分学理论的公理给出。
例如，部分关系是对象上的一个偏序：每个对象都是其自身的一部分（尽管是一个\emph{非真}部分）；没有两个不同的对象可以互为部分；一个对象的部分的部分本身也是该对象的部分。
注意，在这种意义上，“是……的一部分”类似于“是……的子集”，而不类似于“是……的元素”，因为后者既非自反又非传递。
\begin{align*}
& \lforall[x][\Atom{\Obj P}{x,x}] \\
& \lforall[x][\lforall[y][((\Part{x}{y} \land \Part{y}{x})
      \lif \eq[x][y])]] \\
& \lforall[x][\lforall[y][\lforall[z][((\Part{x}{y} \land
        \Part{y}{z}) \lif \Part{x}{z})]]]\\
\intertext{而且，任意两个对象都有一个部分学和（即一个以这两个对象为部分，且在此意义上极小的对象）。}  &
\lforall[x][\lforall[y][\lexists[z][\lforall[u][(\Part{z}{u} \liff
        (\Part{x}{u} \land \Part{y}{u}))]]]]
\end{align*}
这些只是形而上学家所考虑的部分关系的基本原理中的一部分。
然而，若不先引入一些定义的关系，进一步的原理很快就会变得难以表述或书写。
例如，大多数对部分学感兴趣的形而上学家也认为以下是一个有效的原理：每当对象~$x$ 有一个真部分~$y$ 时，它也有一个部分~$z$，该部分与~$y$ 无共同部分，且 $y$ 与 $z$ 的融合即为 $x$。
\end{ex}

\end{document}